\section{One partition for every source set and length}\label{sec:partitions}

The tower in \cref{thm:construction} already uses one syndrome map on
the full binary cube at each level. A general repetition argument fills
the gaps between its lengths. The leftover coordinates are retained and
revealed as raw bits; their cost is included in the label budget.

\begin{theorem}[Theorem F: all-length common partitions]\label{thm:all-length}
\formal{TapeStorageRates/AllLengthConstruction.lean}{TapeStorageRates.AllLengthPartitionTarget,TapeStorageRates.BinaryPartitionFamily.all_lengths,TapeStorageRates.OddGirthLinearSeed.all_length_partitions}
Assume the seed and constants of \cref{thm:construction}, fix
$\Delta\ge D+1$, and put $\varepsilon=C\Delta^{-\alpha}$.
For each positive integer $L$, there are a helper relation $N_L$ on
all $L$ coordinates and a map $f_L:\{0,1\}^L\to Y_L$ with the
following properties. Every incidence degree is at most $\Delta$;
for each fixed radius $r$, all paired distances exceed $r$ for all
sufficiently large $L$. Every nonempty set
\[
 G_z=\{x\in G:f_L(x)=z\},\qquad G\subseteq\{0,1\}^L,
\]
is a storage code for $N_L$. For all sufficiently large $L$,
\begin{equation}\label{eq:all-length-labels}
 |f_L(\{0,1\}^L)|\le 2^{\varepsilon L+\sqrt L}.
\end{equation}
At each such length, simultaneously for every nonempty $G$ and every
real $u$, writing $r(G)=\log_2|G|/L$, one has
\begin{equation}\label{eq:all-length-tail}
 \Pr_{x\ \mathrm{uniform\ on}\ G}
 \left[r(G_{f_L(x)})<r(G)-\varepsilon-L^{-1/2}-u\right]
 \le 2^{-uL}.
\end{equation}
The relations and encoders are chosen before $G$ and $u$.
\end{theorem}
\begin{outline}
More generally, take any positive tower lengths $n_j$ and common
partitions with at most $2^{\varepsilon n_j}$ labels,
$\varepsilon\ge0$, bounded incidence degree and diverging paired
distance. For large $L$, choose
\[
 j(L)=\max\{j\le L:n_j\le\lfloor\sqrt L\rfloor\}.
\]
This index tends to infinity: any fixed $j$ is eligible once
$L\ge\max(j,n_j^2)$. No monotonicity or bound on the growth of $n_j$
is needed. The Lean construction defines the finitely many small
lengths as well, using index zero when the displayed set is empty.

Write $L=cn_{j(L)}+a$, $0\le a<n_{j(L)}$. Apply the tower encoder
independently to $c$ disjoint blocks, and include the remaining $a$
bits in the label. The helper relation is the disjoint union of the
block relations, with empty helper sets at the leftover coordinates.
In each fibre those leftover bits are fixed, so recovery holds there
too. The block degrees are unchanged and leftover ports are isolated.
Since $j(L)\to\infty$, the paired distances diverge. The number of
labels is at most
\[
 (2^{\varepsilon n_{j(L)}})^c2^a
   \le2^{\varepsilon L+a}\le2^{\varepsilon L+\sqrt L}.
\]
The generic index choice is \lean{LeanMachine.SlowIndex.tendsto_index}.
The exact repetition count is
\lean{TapeStorageRates.PartitionPadding.encoder_image_card_le}.
The finite-partition bound
\lean{LeanMachine.FinitePartition.normalized_tail} applies to every
source subset using this same encoder and proves
\cref{eq:all-length-tail}.
\end{outline}

For each fixed cap, the extra rate loss $L^{-1/2}$ tends to zero.
In particular the total exponent in \cref{eq:all-length-labels} is
at most $2\varepsilon L$ eventually; this comparison is checked by
\lean{TapeStorageRates.AllLengthPartitions.encoder_image_card_le_double}.
The theorem concerns uniform input mass on each chosen subset, which
may be correlated or chosen after the encoder. It makes no assertion
about an arbitrary nonuniform probability distribution on that subset.

\begin{corollary}[The unrestricted gap test has a false premise]\label{cor:abstract-gap}
\formal{TapeStorageRates/ConstructiveConsequences.lean}{TapeStorageRates.OddGirthLinearSeed.uniform_familyRateGap_bound,TapeStorageRates.OddGirthLinearSeed.rate_witnesses,TapeStorageRates.OddGirthLinearSeed.not_unbounded_familyRateGap,TapeStorageRates.OddGirthLinearSeed.not_all_powers_familyRateGap}
Assume an odd-characteristic seed with $1/2<\rho<1$.
There are constants $A,\beta>0$, fixed before the function $g$, such
that the following holds: if, for all sufficiently large $\Delta$,
$g(\Delta)\ge0$ and every admissible family at cap $\Delta$ has
limiting upper rate at most $1-g(\Delta)$, then
\[
 g(\Delta)\le A\Delta^{-\beta}
 \qquad\text{for all sufficiently large }\Delta.
\]
For every fixed tape count $k$, the abstract rate-witness conclusion
of Theorem~A holds without a machine assumption. In particular, no
sequence of caps and universal limiting-lower-rate bounds can satisfy
the hypothesis of \cref{cor:gap}.
\end{corollary}
\begin{outline}
At each sufficiently large cap, compare the constructed family's
limiting lower rate with the proposed limiting upper bound. The
constants in the construction do not depend on $g$. For the second
claim, a negative power of $\Delta$ is eventually smaller than
$f_{4k^2+1}(\Delta)$, up to a fixed positive constant. The formal
comparison first uses
\lean{TapeStorageRates.polynomial_mul_loglog_tendsto_zero}, then the
monotonicity of the deficit scale in the tape count.
The stronger sparse-cap premise of \cref{cor:gap} is excluded by
combining the constructed witnesses with
\lean{TapeStorageRates.no_witnesses_of_sparse_rate_bounds}.
\end{outline}

Together with the written seed proof, this closes the unrestricted
rate-gap route at the scale used by Theorem~A. It does not settle (T).
These complete partitions have high-rate typical fibres on all input
coordinates, bounded incidence degree and diverging paired distance.
The construction supplies no realization as cone partitions of a tape
computation. The cost of generating, moving and using the coded
information is absent from these abstract properties.
